\section{Velocity-limited inertia with fixed prototype scores}
\label{sec:inertia}

The preceding experiment transported each query by the first-order field $F$.
We now give the query momentum while holding $F$, the fitted prototypes, and
the class scores fixed. This yields a matched test of whether the law relating
momentum to velocity matters for class assignments outside the observed cores.
Fran\c{c}a et al. used relativistic kinetic energy to design an optimization
algorithm with bounded updates \cite{franca}. Here the trajectory serves as
an inference rule and is assessed against a hard velocity cap and classical
inertia as well as the static readout. Competing fitted potentials also appear
in Galactic Margin Machines, which compare the fields at the query itself
\cite{behnood}.

\subsection{Matched dynamics and a shared preservation bound}

Write $U(x)=-\beta\log\sum_j S_j(x)$, so $F=-\nabla U$. From $x(0)=x$ and
$p(0)=0$, evolve
\begin{equation}
  \dot x=V(p),\qquad \dot p=F(x)-\gamma p,\qquad \gamma=1.4 .
  \label{eq:inertial-state}
\end{equation}
For $r=\norm{p}$ and speed cap $c>0$, the three velocity laws are
\begin{equation}
  V_{\rm classical}(p)=p,\qquad
  V_{\rm hard}(p)=\frac{c}{\max(c,r)}p,\qquad
  V_{\rm rel}(p)=\frac{p}{\sqrt{1+r^2/c^2}} .
  \label{eq:inertial-velocities}
\end{equation}
The momenta are evolved under the same equation in all cases. The hard cap
applies to the velocity in $\dot x$. The relativistic law has unit mass and is
the momentum gradient of $K_{\rm rel}(p)=c^2(\sqrt{1+r^2/c^2}-1)$. The other
velocity laws are also gradients of radial kinetic functions. Classical
inertia uses $r^2/2$, and hard clipping uses $r^2/2$ for $r\le c$ and
$cr-c^2/2$ for $r>c$. For all three, the chain rule gives
$\dd(U+K)/\dd t=-\gamma p^\top V(p)\le 0$. At horizon $t$ the label is
$C(x(t))$, compared to $C(x)$ at initialization. The parameter $c$ is a
speed in standardized feature coordinates per unit of solver time.

Both capped laws obey $\norm{V(p)}\le c$. Because
$\nabla m_{ab}=\mu_a-\mu_b$ for the shared isotropic variance, its norm
is bounded by the largest distance $D_{ab}$ between prototypes in classes
$a$ and $b$. Thus throughout $[0,t]$ every margin for the initial winner
$a$ satisfies
\begin{equation}
  m_{ab}(x(s))\ge m_{ab}(x)-D_{ab}ct,\qquad 0\le s\le t .
  \label{eq:capped-preservation}
\end{equation}
If these lower bounds are strictly positive for every $b\ne a$, the initial
class is preserved throughout the interval in exact arithmetic. The same
test applies to both capped dynamics. At each fixed nonzero momentum the
relativistic law gives a lower speed than the hard cap. A common maximum
speed therefore permits different distances traveled by the two methods.

For smooth $F$ and a fixed query, all three position paths start as
$x(t)=x+\frac12t^2F(x)+O(t^3)$. The hard-capped and classical paths coincide
while the classical momentum stays within the cap. Expanding
$V_{\rm rel}(p)=p-r^2p/(2c^2)+O(r^5)$ and using $p(t)=tF(x)+O(t^2)$ gives
\begin{equation}
  x_{\rm rel}(t)-x_{\rm classical}(t)
    =-\frac{t^4}{8c^2}\norm{F(x)}^2F(x)+O(t^5).
  \label{eq:relativistic-short-time}
\end{equation}
Indeed, the cubic velocity correction first changes displacement at order
$t^4$. The two momentum trajectories then differ by order $t^5$, since
their forcing differs only through the position. The leading displacement
that drives the short-time boundary flux is shared by all three laws.

\subsection{Evaluation on the fixed benchmark roster}

The inertia study uses the synthetic families, gaps, PenDigits and Seeds
cohorts, and core-only feature scaling and prototype fits described above.
The primary bandwidth ratio, horizon and cap ratio are $0.30$, $0.5$ and
$0.15$, respectively. Both capped methods use $c$ equal to the cap ratio
times the fitted prototype diameter. Sensitivity settings cross bandwidth
ratios $0.30$ and $0.45$ with horizons $0.15$, $0.5$ and $1.25$ at the
primary cap, and use cap ratios $0.075$ and $0.30$ at the primary bandwidth
and horizon. Fourth-order Runge--Kutta uses maximum step $0.05$, with a
separate run at $0.025$ to check numerical label sensitivity. Each matched
batch has four single-thread timing repeats with rotated method order.
Timing includes integration where applicable and the final score readout.

The primary comparison is the balanced-accuracy difference between smooth
relativistic and hard-capped inertia. Stage accuracy is the mean of
task-unit balanced accuracies. We also count label changes and their
correct-to-wrong and wrong-to-correct components relative to the static
readout. Normalized endpoint displacement averages the Euclidean query
displacement divided by prototype diameter within each unit and then
across units. Percentile intervals use 10,000 paired seed-block resamples,
preserving comparisons on shared queries and the dependence induced by
repeated gaps. The experimental decision criteria specify a gain of at
least two percentage points over hard clipping with a positive interval
lower endpoint, a positive lower endpoint against static classification,
a maximum median latency ratio of $1.25$ to hard clipping, and at most
$0.1\%$ disagreement under step refinement in every setting. Public
summaries are conditional on the available cohorts. The inertia comparison
reuses the previous benchmark roster and provides a within-protocol
comparison on the same evaluation problems.

\subsection{Inertial transport results}

On the synthetic stage, the primary relativistic-minus-hard difference was
$-0.033$ percentage points (95\% seed-block interval $[-0.141,0.068]$).
The relativistic-minus-static difference was $+0.390$ points
($[-0.158,0.930]$). The capped methods differed on 201 of 51,840 query
evaluations. Mean displacement relative to prototype diameter was $0.0267$
with relativistic inertia and $0.0316$ with hard clipping. The relativistic
law changed 1,171 static labels, including 426 correct-to-wrong and 628
wrong-to-correct transitions. Hard clipping changed 1,371, with 504 and
723 transitions, respectively. The soft cap reduced both categories of
change on these inputs.

On PenDigits and Seeds, the relativistic-minus-hard difference was
$+0.056$ points ($[0,0.167]$) and the relativistic-minus-static difference
was $+0.005$ points ($[-0.222,0.273]$). The capped methods differed on two
of 1,889 query evaluations. The soft and hard mean normalized displacements
were $0.0258$ and $0.0295$. Both stages fell short of the accuracy criteria.
Across all eight settings the largest relativistic gain over hard clipping
was $0.20$ points on the synthetic stage and $0.50$ points on the public
stage. At the longest horizon, relativistic accuracy on the public stage
fell below the static readout by $1.39$ and $1.65$ points at the two bandwidths.
Tables~\ref{tab:inertial-accuracy} and~\ref{tab:inertial-timing} give the
primary accuracy and inference profiles.

\begin{table}[htbp]
\centering
\small
\caption{Unit-macro balanced accuracy (\%) at the primary setting with
fixed Gaussian scores. The static readout evaluates the original query.}
\label{tab:inertial-accuracy}
\InertiaAccuracyTable
\end{table}

\begin{table}[htbp]
\centering
\small
\caption{Single-thread timing and label counts at the primary setting.
Each unit's median batch latency is normalized by query count or divided
by the paired static batch latency before taking the stage median.
Halving the maximum step changed zero labels across all eight settings.
Flips count disagreements with the static readout.}
\label{tab:inertial-timing}
\InertiaTimingTable
\end{table}

The median paired latency ratios for relativistic to hard-capped inference
were $1.019$ and $0.999$, within the threshold on both stages. At the primary
setting each transported readout took 40 field evaluations per query. Median
relativistic-to-static batch ratios were $56.8$ and $55.1$ on this machine.
